\documentclass[aspectratio=169,10pt]{beamer}

\usepackage[utf8]{inputenc}
\usepackage[T1]{fontenc}
\usepackage[french]{babel}
\usepackage{lmodern}
\usepackage{booktabs}
\usepackage{tikz}
\usetikzlibrary{positioning}
\usepackage{pgfplots}
\pgfplotsset{compat=1.18}
\usepackage{hyperref}

\definecolor{navy}{HTML}{12304A}
\definecolor{teal}{HTML}{139A91}
\definecolor{gold}{HTML}{E8A838}
\definecolor{pale}{HTML}{EFF5F8}
\definecolor{muted}{HTML}{657786}

\usetheme{Madrid}
\usecolortheme{default}
\setbeamercolor{structure}{fg=teal}
\setbeamercolor{title}{fg=navy}
\setbeamercolor{subtitle}{fg=teal}
\setbeamercolor{frametitle}{fg=navy,bg=pale}
\setbeamercolor{palette primary}{bg=navy,fg=white}
\setbeamercolor{palette secondary}{bg=navy!86,fg=white}
\setbeamercolor{palette tertiary}{bg=navy!72,fg=white}
\setbeamercolor{palette quaternary}{bg=teal,fg=white}
\setbeamercolor{block title}{bg=navy,fg=white}
\setbeamercolor{block body}{bg=pale,fg=navy}
\setbeamertemplate{navigation symbols}{}
\setbeamertemplate{footline}[frame number]
\setbeamertemplate{itemize item}{\color{teal}\small$\blacktriangleright$}
\setbeamertemplate{itemize subitem}{\color{gold}\scriptsize$\blacktriangleright$}

\title[Sentiment efficace]{Sentiment financier par entropie et compression}
\subtitle{Comparer une représentation compacte à FinBERT}
\author{Votre nom}
\institute{Votre établissement}
\date{Compte-rendu expérimental}

\begin{document}

\begin{frame}[plain]
  \titlepage
  \vspace{-0.5cm}
  \begin{center}
    \small\color{navy!75}
    Objectif : réduire le coût de classification sans masquer le compromis de qualité.
  \end{center}
\end{frame}

\begin{frame}{Question et idée du projet}
  \begin{columns}[T,onlytextwidth]
    \column{0.56\textwidth}
      \begin{block}{Question}
        Peut-on classer le sentiment financier sans dépendre entièrement
        d'embeddings et de Transformers lourds ?
      \end{block}
      \begin{itemize}
        \item Représenter les textes par des mots présents/absents.
        \item Pondérer les mots selon leur gain d'information.
        \item Explorer le principe MDL : choisir l'explication qui minimise
        la description du modèle et du texte restant.
      \end{itemize}
    \column{0.39\textwidth}
      \begin{tikzpicture}[node distance=0.36cm]
        \node[draw=teal,rounded corners,fill=pale,align=center,
          text width=3.45cm,minimum height=0.9cm] (text) {Phrase financière};
        \node[draw=teal,rounded corners,fill=pale,align=center,
          text width=3.45cm,below=of text,minimum height=0.9cm] (info)
          {Mots et information};
        \node[draw=teal,rounded corners,fill=pale,align=center,
          text width=3.45cm,below=of info,minimum height=0.9cm] (class)
          {Sentiment prédit};
        \draw[->,thick,navy] (text) -- (info);
        \draw[->,thick,navy] (info) -- (class);
      \end{tikzpicture}
      \vspace{0.35cm}
      {\scriptsize\color{muted}Deux pistes légères : entropie pondérée et
      classification par compression zlib.}
  \end{columns}
\end{frame}

\begin{frame}{Ce qui a été construit}
  \begin{columns}[T,onlytextwidth]
    \column{0.32\textwidth}
      \begin{block}{1. Entropie}
        Bernoulli Naive Bayes, sélection du vocabulaire par gain d'information,
        jusqu'à 5\,000 caractéristiques.
      \end{block}
    \column{0.32\textwidth}
      \begin{block}{2. MDL / zlib}
        Un dictionnaire par classe, appris sur le train ; prédiction par coût
        du dictionnaire, résidu compressé et a priori.
      \end{block}
    \column{0.32\textwidth}
      \begin{block}{3. Comparaison}
        Baseline majorité, mêmes exemples de test, FinBERT préentraîné en
        inférence seule et rapports réutilisables.
      \end{block}
  \end{columns}
  \vfill
  \begin{alertblock}{Interprétation}
    zlib MDL est une approximation pratique de l'idée MDL du cours, pas un
    code universel optimal. Le modèle entropique n'est pas lui-même un
    algorithme MDL complet.
  \end{alertblock}
\end{frame}

\begin{frame}{Protocole expérimental}
  \begin{columns}[T,onlytextwidth]
    \column{0.48\textwidth}
      \begin{itemize}
        \item Financial PhraseBank, accord des annotateurs $\geq 75\%$.
        \item 3\,448 phrases après déduplication.
        \item Split stratifié : 2\,759 train, 689 test ; graine 42.
        \item 3 classes : négatif, neutre, positif.
        \item FinBERT (\texttt{ProsusAI/finbert}), CPU, 128 tokens, sans
        fine-tuning.
      \end{itemize}
    \column{0.48\textwidth}
      \begin{block}{Pourquoi le F1 macro ?}
        Le test contient 428 exemples neutres, 177 positifs et 84 négatifs.
        Le F1 macro donne le même poids à chaque classe ; l'accuracy seule
        favorise la classe majoritaire.
      \end{block}
      \begin{block}{Artefacts livrés}
        Empreinte du corpus, split et paramètres dans le JSON ; scores par
        classe, matrice de confusion, durées et tailles dans le rapport.
      \end{block}
  \end{columns}
\end{frame}

\begin{frame}{Résultats : qualité sur le test}
  \begin{columns}[T,onlytextwidth]
    \column{0.58\textwidth}
      \centering
      \begin{tikzpicture}
        \begin{axis}[
          xbar,
          width=\textwidth,
          height=5.1cm,
          xmin=0,xmax=1,
          xlabel={Score (0--1)},
          symbolic y coords={FinBERT,Compression,Entropie,Majorité},
          ytick=data,
          y dir=reverse,
          bar width=8pt,
          enlarge y limits=0.2,
          nodes near coords,
          nodes near coords align={horizontal},
          nodes near coords style={font=\scriptsize},
          xtick={0,0.2,0.4,0.6,0.8,1.0},
          grid=major,
          grid style={draw=gray!20},
          tick label style={font=\scriptsize},
          label style={font=\scriptsize},
          axis line style={gray!50},
          tick style={gray!50},
          legend style={at={(0.5,-0.23)},anchor=north,legend columns=2,
            draw=none,font=\scriptsize},
        ]
          \addplot[fill=navy,draw=navy] coordinates {
            (0.2554,Majorité) (0.3206,Entropie) (0.5085,Compression) (0.9308,FinBERT)
          };
          \addlegendentry{F1 macro}
          \addplot[fill=teal,draw=teal] coordinates {
            (0.6212,Majorité) (0.6473,Entropie) (0.7242,Compression) (0.9390,FinBERT)
          };
          \addlegendentry{Accuracy}
        \end{axis}
      \end{tikzpicture}
    \column{0.38\textwidth}
      \begin{alertblock}{À retenir}
        FinBERT obtient le meilleur score. zlib MDL améliore les baselines
        légères, sans atteindre une qualité comparable à FinBERT.
      \end{alertblock}
      {\small
      \begin{tabular}{lrr}
        \toprule
        Modèle & Acc. & F1 macro\\
        \midrule
        Majorité & 0,621 & 0,255\\
        Entropie & 0,647 & 0,321\\
        zlib MDL & 0,724 & 0,509\\
        FinBERT & 0,939 & 0,931\\
        \bottomrule
      \end{tabular}}
  \end{columns}
\end{frame}

\begin{frame}{Compromis mesuré et limites}
  \begin{columns}[T,onlytextwidth]
    \column{0.53\textwidth}
      \begin{block}{Inférence CPU / 1\,000 textes}
        \begin{tabular}{lr}
          \toprule
          zlib MDL & 3,73 s\\
          Entropie pondérée & 6,43 s\\
          FinBERT & 71,76 s\\
          \bottomrule
        \end{tabular}
        \vspace{0.15cm}

        Sur ce lot, zlib est environ 19 fois plus rapide que FinBERT.
      \end{block}
      \begin{block}{Taille enregistrée}
        zlib MDL : 0,099 Mo ; FinBERT : 437,94 Mo.
      \end{block}
    \column{0.43\textwidth}
      \begin{alertblock}{Prudence}
        \begin{itemize}
          \item Un seul split stratifié et un petit corpus de phrases.
          \item Le rappel négatif de zlib n'est que de 9,5\%.
          \item Les tailles ne mesurent pas exactement le même objet :
          modèle sérialisé contre poids seuls de FinBERT.
          \item \texttt{tracemalloc} ne mesure ni toute la RAM native de
          PyTorch ni la RAM pendant l'inférence.
        \end{itemize}
      \end{alertblock}
  \end{columns}
\end{frame}

\begin{frame}{Conclusion et prochaines étapes}
  \begin{block}{Conclusion}
    La compression offre un compromis intéressant en vitesse et taille, mais
    sur ce test elle perd nettement en F1 macro. FinBERT gagne en qualité ;
    zlib MDL reste une piste si les contraintes de ressources dominent.
  \end{block}
  \begin{enumerate}
    \item Répéter avec plusieurs graines et publier moyenne et dispersion.
    \item Tester des textes dédupliqués avec une séparation chronologique.
    \item Améliorer la reconnaissance des sentiments négatifs.
    \item Mesurer la mémoire réelle et garder des tailles de modèles comparables.
  \end{enumerate}
  \vfill
  \begin{center}
    \Large\color{teal}Merci de votre attention
  \end{center}
\end{frame}

\begin{frame}{Références}
  \footnotesize
  \begin{thebibliography}{9}
    \bibitem{malo}
    P. Malo et al. (2014).
    \newblock Good debt or bad debt: Detecting semantic orientations in
    economic texts.
    \newblock \emph{Journal of the Association for Information Science and
    Technology}, 65.
    \newblock Financial PhraseBank, accord de 75\,\%; licence CC BY-NC-SA 3.0.

    \bibitem{araci}
    D. Araci (2019).
    \newblock FinBERT: Financial Sentiment Analysis with Pre-trained Language
    Models.
    \newblock \url{https://arxiv.org/abs/1908.10063}

    \bibitem{rissanen}
    J. Rissanen (1978).
    \newblock Modeling by shortest data description.
    \newblock \emph{Automatica}, 14(5), 465--471.
  \end{thebibliography}
  \vfill
  \scriptsize
  Résultats reproductibles et matrice de confusion :
  \url{https://github.com/EtiennePellissier/Compressed-representation}
\end{frame}

\end{document}
